\documentclass[11pt]{article}
\usepackage[margin=1in]{geometry}
\usepackage{amsmath,amssymb,amsthm}
\usepackage[T1]{fontenc}
\usepackage{lmodern}
\usepackage{xcolor}
\usepackage{booktabs}
\usepackage{array}
\usepackage{longtable}
\usepackage[hidelinks]{hyperref}
\usepackage{url}
\usepackage{newunicodechar}
\newunicodechar{∠}{\ensuremath{\angle}}
\newunicodechar{₁}{\ensuremath{_1}}
\newunicodechar{₂}{\ensuremath{_2}}

\newtheorem{theorem}{Theorem}
\newtheorem{lemma}[theorem]{Lemma}
\theoremstyle{definition}
\newtheorem{definition}{Definition}
\theoremstyle{remark}
\newtheorem{remark}[theorem]{Remark}

\newcommand{\lean}[1]{\texttt{#1}}
% A declaration of one of our developments; every \decl{...} must be in that repository's gate.sh
% REQUIRED list (checked by the name gate).
\DeclareUrlCommand\decl{\urlstyle{tt}\def\UrlBreaks{\do\_\do\.}}
% A Mathlib declaration (pinned revision 0df444a3); not one of ours.
\DeclareUrlCommand\mdecl{\urlstyle{tt}\def\UrlBreaks{\do\_\do\.}}
\newcommand{\RR}{\mathbb{R}}
\newcommand{\PH}[1]{\textcolor{red}{\textbf{[PLACEHOLDER:} #1\textbf{]}}}

% ---------------------------------------------------------------------------------------------
% PINS. The published commits of the three repositories; \src links each File:line to them.
% ---------------------------------------------------------------------------------------------
\expandafter\def\csname commitof:paperh_euler\endcsname{0c95102c792b8d47d2720f4e9fbc39aa5f21a424}
\expandafter\def\csname commitof:paperh_mordell\endcsname{7b2270146195cf2a4c589e403143dba0abd0f6a7}
\expandafter\def\csname commitof:paperh_fermat\endcsname{8f680c86b068196a4d27201209c856a5d038aaa4}
\newcommand{\commitof}[1]{\csname commitof:#1\endcsname}
\newcommand{\repo}[1]{https://github.com/jpbald93/#1}
\newcommand{\src}[4]{\ifcsname commitof:#1\endcsname
  \href{\repo{#1}/blob/\commitof{#1}/#2/#3.lean\#L#4}{\texttt{#3}:#4}\else\texttt{#3}:#4\fi}
\title{Classical triangle inequalities and the Fermat--Napoleon configuration in Lean~4}
\author{Josh Bald\\ \small Independent researcher, Ontario, Canada\\ \small\texttt{jpbald93@gmail.com}\\
\small ORCID: \texttt{0009-0002-1317-6489}
\thanks{\copyright~2026 Josh Bald. Licensed under CC BY 4.0 (\url{https://creativecommons.org/licenses/by/4.0/}). The accompanying Lean code is licensed under the Apache License 2.0.}}
\date{30 October 2026}
\emergencystretch=3em
\begin{document}
\sloppy
\maketitle
\begin{abstract}
We give machine-checked proofs, in Lean~4 with Mathlib, of a body of classical inequalities and
extremal theorems for the triangle, together with the Fermat--Napoleon configuration. For an
arbitrary \lean{Affine.Triangle} in a real inner product space of any dimension, using Mathlib's own
circumcentre, circumradius, incentre, inradius, excentres, exradii, centroid and orthocentre, we
prove Euler's relation $OI^2=R^2-2Rr$, Euler's inequality $R\ge 2r$ \emph{with its equality case}
(equality exactly for equilateral triangles) and the excircle relation $OI_a^2=R^2+2Rr_a$; the
formulas $OG^2=R^2-(a^2+b^2+c^2)/9$, $OH^2=9R^2-(a^2+b^2+c^2)$ and $IG^2=(s^2+5r^2-16Rr)/9$;
Gerretsen's inequalities $16Rr-5r^2\le s^2\le 4R^2+4Rr+3r^2$ and Kooi's sharper upper bound;
the Weitzenb\"ock, Hadwiger--Finsler and triangle isoperimetric inequalities with their equality
cases; the Leibniz relation and the minimality of $\sum|XA_i|^2$ at the centroid; and the
Neuberg--Pedoe inequality for two triangles with equality exactly for similar triangles. We
characterise the equality case of Ptolemy's inequality in any real inner product space (concyclic
points in separated order, or collinear points in one of eight orders). For a point $X$ of the
\emph{closed} triangle we prove the Erdős--Mordell inequality, its key lemma and its equality case
(equilateral triangle, $X$ the centre), and, for interior points, Barrow's inequality. If every
angle is less than $2\pi/3$, the Fermat--Torricelli point exists and minimises $XA+XB+XC$ over the
whole space. In the oriented plane we prove that the segments $AA'$, $BB'$, $CC'$ to the apexes of the
outer equilateral triangles have equal lengths and lie on concurrent lines meeting pairwise at angle
$\pi/3$, that if every angle is less than $2\pi/3$ the common point is the Fermat point $F$ and
$AA'=FA+FB+FC$, and Napoleon's theorem with the centroid
statement. All proofs use only Lean's three standard axioms. Several of these results were
formalised before us, and we reprove them independently: Euler's relation, Euler's inequality and the excircle
relation (two Lean developments, one of them an open Mathlib pull request), Ptolemy's inequality (in
Mathlib since 2022), the Erdős--Mordell inequality and its key lemma (two Lean developments, for
interior points in the plane), algebraic forms of Weitzenb\"ock and Hadwiger--Finsler, and Napoleon's
theorem (Isabelle and several Lean developments, including our own earlier one). For the remaining
statements we found no earlier formal proof in the sources we searched, subject to stated gaps.
\end{abstract}

\section{Introduction}\label{sec:intro}
The inequalities between the circumradius $R$, the inradius $r$, the semiperimeter $s$, the area
$\Delta$ and the sides $a,b,c$ of a triangle form a classical and well-catalogued subject; the
monograph of Bottema, \DJ{}or\dj{}evi\'c, Jani\'c, Mitrinovi\'c and Vasi\'c \cite{bottema} collects
several hundred of them. The best-known are Euler's inequality $R\ge 2r$, which follows from his
relation $OI^2=R(R-2r)$ for the distance between the circumcentre and the incentre; Weitzenb\"ock's
inequality $a^2+b^2+c^2\ge 4\sqrt3\,\Delta$ \cite{weitzenboeck} and its strengthening by Finsler
and Hadwiger \cite{finslerhadwiger}; Gerretsen's two-sided bound on $s^2$ in terms of $R$ and $r$;
and the Erdős--Mordell inequality, posed by Erdős in 1935 \cite{erdos1935} and proved by Mordell and
Barrow \cite{mordellbarrow1937}. This paper formalises these and related results in Lean~4
\cite{lean4} with Mathlib \cite{mathlib}, together with the classical configuration of the
equilateral triangles erected on the sides, which links the Fermat--Torricelli point with Napoleon's
theorem.

\paragraph{Prior formal work, up front.} Several of our statements were formalised before us, and we
claim them only as independent reproofs. Euler's relation $OI^2=R(R-2r)$, Euler's inequality
$2r\le R$ and the excircle relation are proved in Lean by Weiyi Wang, with the same Mathlib objects
as ours, in a public repository (January 2026) \cite{wwylele} and in the open Mathlib pull request
\#44230 (opened 26 September 2026) \cite{pr44230}; the relation and the inequality are also in the
\lean{atlas-lean} repository for the plane with explicitly given incircle data \cite{atlaseuler}.
Ptolemy's inequality has been in Mathlib since 2022 \cite{mathlib} (declaration
\mdecl{EuclideanGeometry.mul_dist_le_mul_dist_add_mul_dist}). The Erdős--Mordell inequality, for a
point in the interior of a triangle in the plane, is proved in Lean in \lean{plby/lean-proofs}
(May 2026) \cite{plby} and in \lean{lean-genius} (June 2026) \cite{lgmordell}, the latter with
the key lemma; the equality case appears in \lean{lean-genius} only as an \lean{axiom}
\cite{lg898}. Weitzenb\"ock's and the Hadwiger--Finsler inequalities with their equality cases are
proved in algebraic form (real numbers $a,b,c,T$ satisfying Heron's relation) in Compfiles
\cite{compfiles} and \lean{lean-genius} \cite{lgweitz,lghf}. Napoleon's theorem is formalised in
Isabelle \cite{afpnapoleon}, in several Lean developments \cite{tanrikolu,shadowbench,atlasnapoleon}
and in our own earlier paper \cite{paperA}. Section~\ref{sec:related} gives the details and dates.

For the other statements we found no earlier formal proof in the sources searched
(Section~\ref{sec:related} lists the gaps): the equality case of Euler's inequality; Gerretsen's
and Kooi's inequalities; the formulas for $OG^2$, $OH^2$ and $IG^2$; the Neuberg--Pedoe
inequality and its equality case; the Leibniz minimum and the triangle isoperimetric inequality
with their equality cases; Weitzenb\"ock and Hadwiger--Finsler for the area of an actual triangle;
the proof of the Erdős--Mordell equality case, and the inequality for the closed triangle in any
dimension; Barrow's inequality; the existence and minimality of the Fermat--Torricelli point; the
characterisation of equality in Ptolemy's inequality; and the Fermat--Napoleon lines. None of
this is new mathematics.

\subsection{Contributions}
Unless stated otherwise, the theorems hold for an arbitrary \mdecl{Affine.Triangle} in a real inner
product space of arbitrary dimension and have no hypothesis beyond the triangle.
\begin{itemize}
\item \emph{The Euler family} (Section~\ref{sec:euler}): Euler's relation, the inequality and its
equality case, the excircle relation, and $OG^2$, $OH^2$, $IG^2$
(Theorems~\ref{thm:euler}--\ref{thm:ogh}).
\item \emph{Inequalities in $R$, $r$, $s$} (Section~\ref{sec:rrs}): the fundamental quartic,
Gerretsen and Kooi (Theorems~\ref{thm:quartic}--\ref{thm:kooi}).
\item \emph{Area inequalities and extremal problems} (Section~\ref{sec:area}): Weitzenb\"ock,
Hadwiger--Finsler, the isoperimetric inequality and the Leibniz minimum, each with its equality case
(Theorems~\ref{thm:weitz}--\ref{thm:leibniz}).
\item \emph{Two triangles} (Section~\ref{sec:pedoe}): Neuberg--Pedoe with its equality case
(Theorem~\ref{thm:pedoe}).
\item \emph{Ptolemy} (Section~\ref{sec:ptolemy}): the inequality (a reproof) and the
characterisation of its equality case (Theorems~\ref{thm:ptolemy} and~\ref{thm:ptolemyeq}).
\item \emph{Erdős--Mordell and Barrow} (Section~\ref{sec:mordell}): key lemma, inequality and
equality case for the closed triangle; Barrow for interior points
(Theorems~\ref{thm:key}--\ref{thm:barrow}).
\item \emph{Fermat and Napoleon} (Section~\ref{sec:fermat}): Fermat--Torricelli existence and
minimality; in the oriented plane, the Fermat--Napoleon lines and Napoleon's theorem
(Theorems~\ref{thm:fermatmin}--\ref{thm:napoleon}).
\end{itemize}
The work is in three Lean developments, \lean{paperh\_euler} (library \lean{EulerIneq}),
\lean{paperh\_mordell} (library \lean{ErdosMordell}) and \lean{paperh\_fermat} (library
\lean{FermatNapoleon}); Table~\ref{tab:thms} lists the Lean name and location of every numbered
theorem. The proofs are by polynomial certificates, as in our earlier papers
\cite{paperA,paperG}: sum-of-squares identities valid modulo Heron's relation, and one barycentric
distance formula, used throughout (Section~\ref{sec:setting}).
\section{Setting, conventions and the two basic tools}\label{sec:setting}
\paragraph{Ambient space and triangle.} Throughout, $V$ is a real inner product space, $P$ a metric
space that is a normed additive torsor over $V$, and $t$ a Mathlib \mdecl{Affine.Triangle} in $P$:
three affinely independent points $A=t_0$, $B=t_1$, $C=t_2$. No finite-dimensionality or dimension
hypothesis is made except in Section~\ref{sec:fermat}, where the Napoleon results assume
$\dim V=2$ and fix an orientation. Non-degeneracy of the triangle is part of
\mdecl{Affine.Triangle} and is never a separate hypothesis. We write
\[
a=|BC|,\quad b=|CA|,\quad c=|AB|,\quad s=\tfrac12(a+b+c)
\]
(\lean{dist (t.points 1) (t.points 2)} and so on; $s$ is \decl{EulerIneq.semiperimeter}), and
$O$, $R$, $I$, $r$, $I_i$, $r_i$, $G$, $H$ for Mathlib's \lean{t.circumcenter},
\lean{t.circumradius}, \lean{t.incenter}, \lean{t.inradius}, \lean{t.excenter \{i\}},
\lean{t.exradius \{i\}}, \lean{t.centroid} and \lean{t.orthocenter}. In Mathlib the incentre and
excentres of a simplex are defined by weights and the radii by the distances to the faces; nothing
in our statements refers to our own versions of these objects. ``Equilateral'' is stated as
$|AB|=|BC|\wedge|BC|=|CA|$.

\paragraph{Area.} Besides $s$ and two auxiliary definitions (\lean{gram} and \lean{sideLen}), the only new definition in \lean{EulerIneq} is the area
\[
\Delta=\tfrac12\,|AB|\,|AC|\,\sin\angle BAC,
\]
with Mathlib's unoriented angle (\decl{EulerIneq.area}). Heron's formula is proved as the bridge
\decl{EulerIneq.two_mul_area_sq}: $(2\Delta)^2=\Gamma(a,b,c)$ with
\[
\Gamma(a,b,c)=\tfrac14\bigl(2a^2b^2+2b^2c^2+2c^2a^2-a^4-b^4-c^4\bigr)
\]
(\decl{EulerIneq.gram}), from the law of cosines and $\sin^2+\cos^2=1$; and $\Delta>0$
(\decl{EulerIneq.area_pos}). The other bridges tie Mathlib's objects to the side lengths:
$h_ia_i=2\Delta$ for Mathlib's heights (\decl{EulerIneq.height_mul_sideLen}); $4R\Delta=abc$
(\decl{EulerIneq.four_mul_circumradius_mul_area}, through $4R^2\Gamma=a^2b^2c^2$, adapted from our
earlier Blundon development \cite{paperD}); Mathlib's excentre weights are proportional to
$(\pm a,\pm b,\pm c)$ with the minus sign at the excentre's vertex
(\decl{EulerIneq.excenterWeights_eq}), so that $I$ has weights $(a,b,c)/(2s)$
(\decl{EulerIneq.incenter_weights}); $r=\Delta/s$ (\decl{EulerIneq.inradius_eq},
\decl{EulerIneq.inradius_mul}) and Mathlib's exradius formula (\decl{EulerIneq.exradius_eq}); and
the strict triangle inequalities $-a+b+c>0$ etc., obtained from Mathlib's positivity of the
unnormalised excentre weight sums (\decl{EulerIneq.sum_signed_sideLen_pos},
\decl{EulerIneq.triangle_pos}).

\paragraph{Points in the closed triangle; feet of perpendiculars.} For weights $w_0,w_1,w_2$ with
$w_0+w_1+w_2=1$ we write $w_0A+w_1B+w_2C$ for Mathlib's \mdecl{Finset.affineCombination}. A point
$X$ is \emph{in the closed triangle} if $X=w_0A+w_1B+w_2C$ for some weights with all $w_i\ge0$, and
\emph{strictly inside} if all $w_i>0$; these are the literal hypotheses of
Section~\ref{sec:mordell}. The distance from $X$ to the line $YZ$ is $|X\,\pi_{YZ}(X)|$, where
$\pi_{YZ}$ is Mathlib's \mdecl{EuclideanGeometry.orthogonalProjection} onto
\lean{line[\(\mathbb{R}\), Y, Z]}; we write $d_A(X)$, $d_B(X)$, $d_C(X)$ for the distances to the
lines $BC$, $CA$, $AB$. (The library abbreviates this as \decl{ErdosMordell.footDist}, but the
main statements are written with the projection itself.)

\paragraph{Tool 1: the barycentric distance formula.} The identity behind Euler's relation, the
excircle relation, $OG^2$, $IG^2$ and the Leibniz minimum is the classical formula of Lagrange and
Leibniz for an affine combination, proved once for three points in any inner product space
(\decl{EulerIneq.dist_affineCombination_sq}):
\begin{equation}\label{eq:leibniz}
\Bigl|\sum_i w_iP_i-Q\Bigr|^2=\sum_i w_i|P_iQ|^2-\sum_{i<j}w_iw_j|P_iP_j|^2\qquad
\Bigl(\sum_i w_i=1\Bigr).
\end{equation}
With $Q=O$ the first sum is $R^2$ (\decl{EulerIneq.dist_affineCombination_circumcenter_sq}); for two
affine combinations with weight difference $d=w-v$ the distance squared is
$-\sum_{i<j}d_id_j|P_iP_j|^2$ (\decl{EulerIneq.dist_affineCombination_affineCombination_sq}).
Since Mathlib's $I$, $I_i$ and $G$ are affine combinations with explicit weights, every distance
between these points and $O$ becomes a rational function of $a,b,c$, which is then rewritten in
$R$, $r$, $s$, $\Delta$ by the bridges.

\paragraph{Tool 2: certificates modulo Heron.} Every inequality between $R$, $r$, $s$, $\Delta$ and
the sides is reduced to a polynomial statement in $a,b,c,\Delta$ subject to $(2\Delta)^2=\Gamma$,
$\Delta>0$ and the strict triangle inequalities, and then closed by an explicit identity: a
sum of squares with nonnegative weights, or a product of known-sign factors, found with SymPy and
checked by \lean{linear\_combination}. The equality cases come out of the same identity: equality
forces every square to vanish. As in our earlier papers, the identities were checked numerically
(mpmath, 30 digits, thousands of random, obtuse and near-degenerate triangles) before the Lean
proofs were written; this is evidence, not part of the proofs. We write $x=s-a$, $y=s-b$, $z=s-c$
(the tangent lengths, all positive), so that $a=y+z$ etc.\ (the Ravi substitution).
\section{The Euler family}\label{sec:euler}
\begin{longtable}{@{}l p{3.7cm} p{6.3cm} l@{}}
\caption{Numbered theorems, their Lean names and locations. Namespaces: \lean{EulerIneq}
(development \lean{paperh\_euler}), \lean{ErdosMordell} (\lean{paperh\_mordell}) and
\lean{FermatNapoleon} (\lean{paperh\_fermat}); the namespace is omitted in the table.
``hyp.''\ lists hypotheses beyond the triangle $t$ (or beyond the four points for Ptolemy).
Rows marked R are independent reproofs of results formalised earlier (Section~\ref{sec:related}).}
\label{tab:thms}\\
\toprule
Thm & content & Lean name & File:line\\
\midrule\endfirsthead
\toprule Thm & content & Lean name & File:line\\ \midrule\endhead
\bottomrule\endlastfoot
\ref{thm:euler} R & $OI^2=R^2-2Rr$ & \decl{euler} & \src{paperh_euler}{EulerIneq}{Main}{22}\\
\ref{thm:eulerineq} R & $2r\le R$ & \decl{euler_inequality} & \src{paperh_euler}{EulerIneq}{Main}{73}\\
\ref{thm:eulerineq} & equality iff equilateral & \decl{euler_inequality_eq_iff} & \src{paperh_euler}{EulerIneq}{Main}{91}\\
\ref{thm:eulerex} R & $OI_i^2=R^2+2Rr_i$ & \decl{euler_excenter} & \src{paperh_euler}{EulerIneq}{Main}{117}\\
\ref{thm:ogh} & $OG^2$ & \decl{dist_circumcenter_centroid_sq} & \src{paperh_euler}{EulerIneq}{Main}{349}\\
\ref{thm:ogh} & $OH^2$ & \decl{dist_circumcenter_orthocenter_sq} & \src{paperh_euler}{EulerIneq}{Main}{358}\\
\ref{thm:ogh} & $IG^2$ & \decl{dist_incenter_centroid_sq} & \src{paperh_euler}{EulerIneq}{Main}{370}\\
\midrule
\ref{thm:quartic} & fundamental quartic & \decl{fundamental_quartic} & \src{paperh_euler}{EulerIneq}{Main}{280}\\
\ref{thm:gerretsen} & Gerretsen & \decl{gerretsen} & \src{paperh_euler}{EulerIneq}{Main}{329}\\
\ref{thm:gerretsen} & cross-check & \decl{gerretsen_cross_check} & \src{paperh_euler}{EulerIneq}{Extremal}{153}\\
\ref{thm:kooi} & Kooi & \decl{kooi} & \src{paperh_euler}{EulerIneq}{Kooi}{46}\\
\ref{thm:kooi} & Kooi $\le$ Gerretsen & \decl{kooi_le_gerretsen} & \src{paperh_euler}{EulerIneq}{Kooi}{54}\\
\midrule
\ref{thm:weitz} & Weitzenb\"ock & \decl{weitzenboeck}, \decl{weitzenboeck_eq_iff} & \src{paperh_euler}{EulerIneq}{Main}{236}, 246\\
\ref{thm:hf} & Hadwiger--Finsler & \decl{hadwiger_finsler}, \decl{hadwiger_finsler_eq_iff} & \src{paperh_euler}{EulerIneq}{Main}{183}, 203\\
\ref{thm:iso} & isoperimetric & \decl{isoperimetric}, \decl{isoperimetric_eq_iff} & \src{paperh_euler}{EulerIneq}{Extremal}{65}, 93\\
\ref{thm:leibniz} & Leibniz relation & \decl{leibniz_relation} & \src{paperh_euler}{EulerIneq}{Extremal}{23}\\
\ref{thm:leibniz} & minimum, equality & \decl{leibniz_min}, \decl{leibniz_min_eq_iff} & \src{paperh_euler}{EulerIneq}{Extremal}{35}, 42\\
\midrule
\ref{thm:pedoe} & Neuberg--Pedoe & \decl{pedoe}, \decl{pedoe_eq_iff} & \src{paperh_euler}{EulerIneq}{Pedoe}{43}, 75\\
\midrule
\ref{thm:ptolemy} R & Ptolemy's inequality & \decl{ptolemy_inequality} & \src{paperh_euler}{EulerIneq}{Ptolemy}{20}\\
\ref{thm:ptolemyeq} & equality characterisation; four pairwise distinct points & \decl{ptolemy_eq_iff} & \src{paperh_fermat}{FermatNapoleon}{Ptolemy}{239}\\
\ref{thm:ptolemyeq} & collinear case & \decl{ptolemy_collinear}, \decl{ptolemy_real} & \src{paperh_fermat}{FermatNapoleon}{Ptolemy}{72}, 52\\
\midrule
\ref{thm:key} R & key lemma, three vertices; $X$ in closed $t$ & \decl{key_lemma_zero}, \decl{key_lemma_one}, \decl{key_lemma_two} & \src{paperh_mordell}{ErdosMordell}{Mordell}{130}, 146, 162\\
\ref{thm:em} R & Erdős--Mordell; $X$ in closed $t$ & \decl{erdos_mordell} & \src{paperh_mordell}{ErdosMordell}{Mordell}{179}\\
\ref{thm:em} & equality case; $X$ in closed $t$ & \decl{erdos_mordell_eq_iff} & \src{paperh_mordell}{ErdosMordell}{Mordell}{386}\\
\ref{thm:barrow} & Barrow; $X$ strictly inside, bisector feet & \decl{barrow} & \src{paperh_mordell}{ErdosMordell}{Barrow}{283}\\
\midrule
\ref{thm:fermatmin} & minimality; three $2\pi/3$ angles & \decl{fermat_min} & \src{paperh_mordell}{ErdosMordell}{Fermat}{44}\\
\ref{thm:fermat} & existence; angles $<2\pi/3$ & \decl{fermat_exists}, \decl{fermat_torricelli} & \src{paperh_mordell}{ErdosMordell}{Fermat}{163}, 291\\
\ref{thm:fn} & outward apexes; $\dim 2$, $o$ & \decl{apex_outward} & \src{paperh_fermat}{FermatNapoleon}{Main}{20}\\
\ref{thm:fn} & $AA'=BB'=CC'$; $\dim 2$, $o$ & \decl{napoleon_dist_eq} & \src{paperh_fermat}{FermatNapoleon}{Main}{63}\\
\ref{thm:fn} & concurrence; $\dim 2$, $o$ & \decl{napoleon_concurrent}, \decl{fermat_point_spec} & \src{paperh_fermat}{FermatNapoleon}{Main}{99}, 72\\
\ref{thm:fn} & Fermat point on the lines; $\dim 2$, $o$, angles $<2\pi/3$ & \decl{napoleon_fermat} & \src{paperh_fermat}{FermatNapoleon}{Main}{111}\\
\ref{thm:fn} & angles between the lines; $\dim 2$, $o$ & \decl{napoleon_angle} & \src{paperh_fermat}{FermatNapoleon}{Main}{258}\\
\ref{thm:napoleon} R & Napoleon; $\dim 2$, $o$ & \decl{napoleon} & \src{paperh_fermat}{FermatNapoleon}{Main}{222}\\
\end{longtable}

The theorems of Sections~\ref{sec:euler}--\ref{sec:pedoe} are in \lean{paperh\_euler}; each has
the triangle (or, for Pedoe, the two triangles) as its only hypothesis. File:line references are
to the library files; the commits are given in Section~\ref{sec:avail}.

\begin{theorem}[Euler's relation; reproof]\label{thm:euler}
For every triangle, $|IO|^2=R^2-2Rr$.
\end{theorem}
\begin{theorem}[Euler's inequality and its equality case]\label{thm:eulerineq}
For every triangle, $2r\le R$; and $2r=R$ if and only if the triangle is equilateral.
\end{theorem}
\begin{theorem}[Euler for the excircles; reproof]\label{thm:eulerex}
For every triangle and every vertex $i$, $|I_iO|^2=R^2+2Rr_i$, where $I_i$ and $r_i$ are Mathlib's
\lean{t.excenter \{i\}} and \lean{t.exradius \{i\}} (the excircle opposite vertex $i$).
\end{theorem}
\begin{theorem}[$OG$, $OH$, $IG$]\label{thm:ogh}
For every triangle,
\[
|OG|^2=R^2-\tfrac19(a^2+b^2+c^2),\qquad |OH|^2=9R^2-(a^2+b^2+c^2),\qquad
|IG|^2=\tfrac19\bigl(s^2+5r^2-16Rr\bigr).
\]
\end{theorem}

\paragraph{Proofs.} By the incentre bridge, $I$ has weights $(a,b,c)/(a+b+c)$, so
\eqref{eq:leibniz} with $Q=O$ gives $|IO|^2=R^2-abc/(a+b+c)$, and $abc/(a+b+c)=2Rr$ by $4R\Delta=abc$
and $r(a+b+c)=2\Delta$. For the excentre opposite $A$ the weights are $(-a,b,c)/(-a+b+c)$ and the
same computation gives $R^2+abc/(-a+b+c)=R^2+2Rr_a$; the other two excentres are analogous. For the
inequality, the key identity (\decl{EulerIneq.four_mul_area_mul_circumradius_sub}) is
\[
4\Delta(R-2r)=\tfrac12\bigl[(-a+b+c)(b-c)^2+(a-b+c)(c-a)^2+(a+b-c)(a-b)^2\bigr],
\]
a polynomial identity once $R=abc/(4\Delta)$ and $r=2\Delta/(a+b+c)$ are substituted. Since
$\Delta>0$ and the three linear factors are positive, the right side is nonnegative, and it vanishes
exactly when $a=b=c$ (the auxiliary lemmas \decl{EulerIneq.sum_three_eq_zero} and
\decl{EulerIneq.eq_zero_of_pos_mul_sq}). Note that Euler's relation does not give the equality case
by itself: $|IO|=0$ means $I=O$, and one still has to show that this forces the triangle to be
equilateral; the identity above does both at once. For $OG^2$ apply \eqref{eq:leibniz} with weights
$\tfrac13$ and $Q=O$. For $OH^2$, Mathlib gives $H=O+3(G-O)$
(\mdecl{Affine.Triangle.orthocenter_eq_smul_vsub_vadd_circumcenter}), so $|OH|^2=9|OG|^2$. For
$IG^2$ use the distance between two affine combinations with $d=(a,b,c)/(2s)-(\tfrac13,\tfrac13,\tfrac13)$;
the result is a symmetric rational function of $a,b,c$, which equals $(s^2+5r^2-16Rr)/9$ after substituting
$r=2\Delta/(a+b+c)$ and $4R\Delta=abc$ and using Heron's relation (one \lean{linear\_combination}).
\section{Inequalities in $R$, $r$ and $s$: the fundamental quartic, Gerretsen and Kooi}\label{sec:rrs}
\begin{theorem}[The fundamental quartic]\label{thm:quartic}
For every triangle,
$s^4-2(2R^2+10Rr-r^2)\,s^2+r(4R+r)^3\le0$.
\end{theorem}
\begin{theorem}[Gerretsen's inequalities]\label{thm:gerretsen}
For every triangle, $16Rr-5r^2\le s^2\le 4R^2+4Rr+3r^2$. Moreover
$(R-r)(R-2r)\ge0$ (\decl{EulerIneq.gerretsen_cross_check}).
\end{theorem}
\begin{theorem}[Kooi's inequality]\label{thm:kooi}
For every triangle,
\[
s^2\le\frac{R(4R+r)^2}{2(2R-r)}\qquad\text{and}\qquad \frac{R(4R+r)^2}{2(2R-r)}\le 4R^2+4Rr+3r^2,
\]
so Kooi's bound is at least as strong as the upper Gerretsen bound.
\end{theorem}

The denominator $2(2R-r)$ is positive by Theorem~\ref{thm:eulerineq}; the statement uses Lean's
real division, but no case relies on the convention $x/0=0$.

\paragraph{The quartic.} The quartic is the discriminant form of Blundon's fundamental inequality
\cite{blundon1965}: in the Ravi variables $x,y,z$ (so $s=x+y+z$), the relations
$r^2s=xyz$ and $4Rrs=(y+z)(z+x)(x+y)$ hold, and the certificate
\decl{EulerIneq.core_identity},
\[
s^3\bigl(4r^2 f(s^2)+\bigl((y-x)(z-y)(x-z)\bigr)^2\bigr)=0,\qquad
f(u)=u^2-2(2R^2+10Rr-r^2)u+r(4R+r)^3,
\]
holds modulo these two relations (the explicit cofactors are in the source). Since $s>0$ and
$r>0$, $f(s^2)\le0$. The certificate and the quartic are copied, with attribution in the source,
from our earlier Blundon development \cite{paperD}; here they are restated with Mathlib's $R$ and
$r$ and our $s$.

\paragraph{Gerretsen.} The roots of $f$ are $2R^2+10Rr-r^2\pm2(R-2r)\sqrt{R(R-2r)}$, and
Gerretsen's bounds are the observation that this interval lies inside
$[16Rr-5r^2,\,4R^2+4Rr+3r^2]$. In Lean this is done without square roots: from $f(s^2)\le 0$ and
$R\ge 2r$ one gets
\[
\bigl(s^2-(2R^2+10Rr-r^2)\bigr)^2\le\bigl(2(R-r)(R-2r)\bigr)^2
\]
by \lean{nlinarith} with the hints $(R-2r)^2\ge0$ and $(R-2r)^2r^2\ge0$, and then
$|s^2-(2R^2+10Rr-r^2)|\le 2(R-r)(R-2r)$, which unfolds to both bounds (equality in either bound holds only for the equilateral triangle, where $R=2r$, and is not formalised). The cross-check states what
Gerretsen's two bounds give when combined: $(R-r)(R-2r)\ge0$. This does \emph{not} imply $R\ge2r$,
so Euler's inequality is not a consequence of Gerretsen's bounds alone (we use it in their proof).

\paragraph{Kooi.} Write $K=R(4R+r)^2/(2(2R-r))$. The algebraic lemma
\decl{EulerIneq.kooi_of_quartic} shows: if $r>0$, $2r\le R$ and $f(u)\le0$, then $u\le K$. Its proof
uses the identity
\[
M^2-\bigl(2(2R-r)(u-B)\bigr)^2=\bigl(r(R-2r)(4R+r)\bigr)^2-4(2R-r)^2f(u),
\]
with $B=2R^2+10Rr-r^2$ and $M=(R-2r)(8R^2-12Rr+r^2)$; here $M\ge0$ because $R\ge2r$ makes both
factors nonnegative. Hence $2(2R-r)(u-B)\le M$, which is $u\le K$. Kooi's inequality is this lemma at
$u=s^2$. The comparison with Gerretsen is
$4R^2+4Rr+3r^2-K=3r^2(R-2r)/(2(2R-r))\ge0$, closed after clearing the denominator by
\lean{nlinarith} with the hint $r^2(R-2r)\ge0$. Numerically, equality in Kooi's inequality holds only for
the equilateral triangle; we have not formalised that equality case. In \lean{kooi\_of\_quartic}, since $(u-B)^2-f(u)=4R(R-2r)^3$, the hypotheses $r>0$ and $f(u)\le0$
together with $R\ge0$ already force $2r\le R$ (the lemma states $2r\le R$ explicitly to avoid a separate
positivity hypothesis on $R$); Theorem~\ref{thm:kooi} itself has no hypothesis beyond the triangle.
\section{Area inequalities and two extremal problems}\label{sec:area}
\begin{theorem}[Weitzenb\"ock]\label{thm:weitz}
For every triangle, $4\sqrt3\,\Delta\le a^2+b^2+c^2$, with equality if and only if the triangle is
equilateral.
\end{theorem}
\begin{theorem}[Hadwiger--Finsler]\label{thm:hf}
For every triangle,
$4\sqrt3\,\Delta+(a-b)^2+(b-c)^2+(c-a)^2\le a^2+b^2+c^2$, with equality if and only if the
triangle is equilateral.
\end{theorem}
\begin{theorem}[Isoperimetric inequality for triangles]\label{thm:iso}
For every triangle, $12\sqrt3\,\Delta\le(a+b+c)^2$, with equality if and only if the triangle is
equilateral.
\end{theorem}
\begin{theorem}[Leibniz relation and minimum]\label{thm:leibniz}
For every triangle and every point $X$ of the space,
\[
|XA|^2+|XB|^2+|XC|^2=3|XG|^2+\tfrac13(a^2+b^2+c^2);
\]
hence $|XA|^2+|XB|^2+|XC|^2\ge\tfrac13(a^2+b^2+c^2)$, with equality if and only if $X=G$.
\end{theorem}
Here $\Delta$ is the area of the triangle as defined in Section~\ref{sec:setting}, and in
Theorem~\ref{thm:leibniz} $X$ ranges over all of $P$, not only the plane of the triangle.

\paragraph{Proofs.} With $L=2(ab+bc+ca)-(a^2+b^2+c^2)$, the Hadwiger--Finsler inequality says
$4\sqrt3\,\Delta\le L$. Heron's relation gives the identity (\decl{EulerIneq.hf_identity})
\[
L^2-(4\sqrt3\,\Delta)^2=2\bigl[\bigl((a-b+c)(c-a)\bigr)^2+\bigl((a+b-c)(a-b)\bigr)^2+\bigl((-a+b+c)(b-c)\bigr)^2\bigr],
\]
and $L=4(xy+yz+zx)>0$ in the Ravi variables (\decl{EulerIneq.hf_L_nonneg}); taking square roots of
nonnegative quantities gives the inequality. Equality forces the three squares to vanish, and since
the linear factors are positive, $a=b=c$. Weitzenb\"ock follows by adding
$(a-b)^2+(b-c)^2+(c-a)^2\ge0$; its equality case uses both that sum and the Hadwiger--Finsler
equality case. For the isoperimetric inequality, Heron gives
$(12\sqrt3\,\Delta)^2=27\,p\,(-a+b+c)(a-b+c)(a+b-c)$ with $p=a+b+c$, while
$p^4=p\,\bigl((-a+b+c)+(a-b+c)+(a+b-c)\bigr)^3$; three-term AM--GM in the explicit form
(\decl{EulerIneq.amgm_three})
\[
(x+y+z)^3-27xyz=\tfrac12(x+y+z)\bigl[(x-y)^2+(y-z)^2+(z-x)^2\bigr]+3\bigl[x(y-z)^2+y(z-x)^2+z(x-y)^2\bigr]
\]
finishes the inequality, and the same identity gives the equality case. The Leibniz relation is
\eqref{eq:leibniz} with weights $\tfrac13$ and $Q=X$; the minimum and its equality case follow
because $|XG|^2\ge0$ with equality iff $X=G$.

\section{Two triangles: the Neuberg--Pedoe inequality}\label{sec:pedoe}
\begin{theorem}[Neuberg--Pedoe]\label{thm:pedoe}
Let $t$ and $t'$ be triangles, possibly in different real inner product spaces, with sides
$a,b,c$ and $a',b',c'$ (in the same vertex labelling) and areas $\Delta$, $\Delta'$. Then
\[
a'^2(b^2+c^2-a^2)+b'^2(c^2+a^2-b^2)+c'^2(a^2+b^2-c^2)\ \ge\ 16\,\Delta\Delta',
\]
with equality if and only if there is $k>0$ with $a'=ka$, $b'=kb$ and $c'=kc$.
\end{theorem}
The equality condition is similarity with the given correspondence of vertices; the statement is
not symmetric under relabelling $t'$ alone.

\paragraph{Proof.} Write $S=a^2+b^2+c^2$, $T=a'^2+b'^2+c'^2$ and $\Lambda$ for the left side. The
certificate \decl{EulerIneq.pedoe_identity} is
\[
2ST\bigl(\Lambda-16\Delta\Delta'\bigr)=\bigl(4T\Delta-4S\Delta'\bigr)^2
+2\bigl[(Ta^2-Sa'^2)^2+(Tb^2-Sb'^2)^2+(Tc^2-Sc'^2)^2\bigr],
\]
valid for real numbers whenever $(2\Delta)^2=\Gamma(a,b,c)$ and $(2\Delta')^2=\Gamma(a',b',c')$.
Since $S,T>0$, the inequality follows. If equality holds, all the squares vanish, so
$a'^2=(T/S)a^2$ and so on, and $k=\sqrt{T/S}$; conversely, for similar triangles $\Delta'=k^2\Delta$ and
both sides equal $16k^2\Delta^2$. No Lorentzian Cauchy--Schwarz step is needed. Pedoe's notes
\cite{pedoe1941,pedoe1942} give the inequality and its equality case; the name Neuberg--Pedoe follows
common usage, and we have not checked Neuberg's original nineteenth-century source.
\section{Ptolemy's inequality and its equality case}\label{sec:ptolemy}
\begin{theorem}[Ptolemy's inequality; reproof]\label{thm:ptolemy}
For any four points $A,B,C,D$ of $P$, $|AC|\,|BD|\le|AB|\,|CD|+|BC|\,|AD|$.
\end{theorem}
\begin{theorem}[Equality in Ptolemy's inequality]\label{thm:ptolemyeq}
Let $A,B,C,D$ be four pairwise distinct points of $P$ (six hypotheses $A\ne B,\dots,C\ne D$).
Then $|AC|\,|BD|=|AB|\,|CD|+|BC|\,|AD|$ if and only if one of the following holds:
\begin{enumerate}
\item[(a)] $A,B,C,D$ are cospherical (Mathlib's \mdecl{EuclideanGeometry.Cospherical}: they lie on
a common sphere of $P$) and there is a point $X$ strictly between $A$ and $C$ and strictly between
$B$ and $D$ (the open diagonals meet); or
\item[(b)] $A,B,C,D$ are collinear, $B$ or $D$ is strictly between $A$ and $C$, and $A$ or $C$ is
strictly between $B$ and $D$.
\end{enumerate}
\end{theorem}
Theorem~\ref{thm:ptolemy} is Mathlib's \mdecl{EuclideanGeometry.mul_dist_le_mul_dist_add_mul_dist}
(Kudryashov, 2022), restated in our namespace and proved by that lemma; we list it only for
completeness. In Theorem~\ref{thm:ptolemyeq}, ``strictly between'' is Mathlib's \lean{Sbtw}; case (b)
is exactly the eight cyclic orders $ABCD$, $BCDA$, $CDAB$, $DABC$ and their reversals on the line,
which the developers checked against all 24 orders before proving it. No dimension hypothesis is
needed: in a space of dimension $>2$, ``cospherical'' for four non-collinear points with crossing
diagonals means that they lie on a circle in their common plane. The distinctness hypotheses are
needed: for $A=B$ and $C,D$ not collinear with $A$ the equality holds trivially but neither (a) nor (b)
does, and the tampering suite checks that removing $A\ne B$ breaks the build.

\paragraph{Proof.} The direction ``(a) $\Rightarrow$ equality'' is Mathlib's
\mdecl{EuclideanGeometry.mul_dist_add_mul_dist_eq_mul_dist_of_cospherical}, applied with the angle
$\angle AXC=\pi$ supplied by the strict betweenness. For (b) the points are transported to the real
line along the line $AC$, where the statement becomes a statement about four distinct reals,
$|a-c||b-d|=|a-b||c-d|+|b-c||a-d|$, decided by a case split on the $2^6$ order patterns
(\decl{FermatNapoleon.ptolemy_real}, \decl{FermatNapoleon.ptolemy_collinear}). For the converse,
put $b=B-A$, $c=C-A$, $d=D-A$ and invert about $A$: $x\mapsto x'=x/|x|^2$. Since
$|x'-y'|=|x-y|/(|x||y|)$, the equality becomes $|b'-c'|+|c'-d'|=|b'-d'|$, so $c'=(1-\mu)b'+\mu d'$
with $0\le\mu\le1$ (\lean{ptolemy\_vec}). If $b'$ and $d'$ are parallel, all four points are
collinear and the real-line case applies. Otherwise there is a vector $m$ with
$\langle m,b'\rangle=\langle m,d'\rangle=\tfrac12$, hence $\langle m,c'\rangle=\tfrac12$; the sphere
with centre $A+m$ through $A$ then contains $B$, $C$ and $D$ (because $\langle m,x'\rangle=\tfrac12$
is equivalent to $|x-m|=|m|$). Finally, with $\kappa=(1-\mu)|b'|^2+\mu|d'|^2$, the point
$X=A+c'/\kappa$ lies on $AC$ with parameter $(|c|^2\kappa)^{-1}\in(0,1)$ and on $BD$ with parameter
$\mu/(\kappa|d|^2)\in(0,1)$, which gives the crossing diagonals. The inversion is done by hand on vectors.
\section{The Erdős--Mordell inequality, its equality case, and Barrow's inequality}\label{sec:mordell}
The theorems of this section are in \lean{paperh\_mordell}. Recall (Section~\ref{sec:setting}) that
``$X$ in the closed triangle'' is the hypothesis $\exists w,\ (\forall i,\ w_i\ge0)\wedge\sum_iw_i=1
\wedge X=w_0A+w_1B+w_2C$, and that $d_A(X)$, $d_B(X)$, $d_C(X)$ are the distances from $X$ to its
orthogonal projections onto the lines $BC$, $CA$, $AB$.

\begin{theorem}[Key lemma; reproof, extended]\label{thm:key}
For every triangle and every $X$ in the closed triangle,
\[
b\,d_C(X)+c\,d_B(X)\le a\,|XA|,\qquad c\,d_A(X)+a\,d_C(X)\le b\,|XB|,\qquad
a\,d_B(X)+b\,d_A(X)\le c\,|XC|.
\]
\end{theorem}
\begin{theorem}[Erdős--Mordell and its equality case]\label{thm:em}
For every triangle and every $X$ in the closed triangle,
\[
2\bigl(d_A(X)+d_B(X)+d_C(X)\bigr)\le|XA|+|XB|+|XC|,
\]
with equality if and only if the triangle is equilateral and $X=G$ (Mathlib's
\lean{Finset.univ.centroid \(\mathbb{R}\) t.points}).
\end{theorem}
\begin{theorem}[Barrow's inequality]\label{thm:barrow}
Let $X$ be strictly inside the triangle (all three weights positive), and let $U$, $V'$, $W$ be points
with $U$ weakly between $B$ and $C$, $V'$ weakly between $C$ and $A$, $W$ weakly between $A$ and $B$
(Mathlib's \lean{Wbtw}), such that $\angle BXU=\angle UXC$, $\angle CXV'=\angle V'XA$ and
$\angle AXW=\angle WXB$. Then
\[
2\bigl(|XU|+|XV'|+|XW|\bigr)\le|XA|+|XB|+|XC|.
\]
\end{theorem}
The pairing in the key lemma matters for the proof: the unpaired form $b\,d_B+c\,d_C\le a|XA|$
is also true (it is $|XA|+d_A\ge h_a$), but summing it does not give the Erdős--Mordell inequality;
the reflected pairing $b\,d_C+c\,d_B\le a|XA|$ is the one needed.
Compared with the earlier Lean proofs \cite{plby,lgmordell}, which assume a two-dimensional space and
$X$ in the topological interior of the convex hull, Theorems~\ref{thm:key} and~\ref{thm:em} allow
any dimension and the closed triangle, including the boundary and the vertices; this is a mild
generalisation, not new mathematics. In Barrow's inequality, $U$ is the foot on $BC$ of the internal
bisector of $\angle BXC$; we require $X$ strictly inside because on the boundary $X$ can coincide
with a vertex and the bisector hypotheses degenerate. Barrow's inequality strengthens Erdős--Mordell,
since $|XU|\ge d_A(X)$; we do not derive one from the other in Lean.

\paragraph{Key lemma.} At vertex $A$ write $X-A=w_1(B-A)+w_2(C-A)$ with $w_1,w_2\ge0$, let
$k=\langle B-A,C-A\rangle$ and $\Gamma=b^2c^2-k^2>0$ (the Gram determinant). Then
$c\,d_C=w_2\sqrt\Gamma$ and $b\,d_B=w_1\sqrt\Gamma$ (from the explicit form of Mathlib's projection
onto a line, \decl{ErdosMordell.proj_eq_lineMap}), and $|XA|^2=w_1^2c^2+2w_1w_2k+w_2^2b^2$. After
multiplying by $bc$ and squaring, the key lemma is the identity (\lean{key\_scalar\_sq})
\[
\bigl(bc\cdot a|XA|\bigr)^2-\bigl(bc(b\,d_C+c\,d_B)\bigr)^2
=\bigl(b^2c^2w_1+b^2kw_2-c^2kw_1-b^2c^2w_2\bigr)^2,
\]
which uses $a^2=b^2+c^2-2k$. The key lemma itself is classical (see Kazarinoff
\cite{kazarinoff1957}); our proof of it is this coordinate identity. The other two vertices are relabellings.

\paragraph{Erdős--Mordell.} Multiply the three key-lemma inequalities by $bc$, $ca$, $ab$, add, and
divide by $abc$; what remains is
$\sum_{\mathrm{cyc}}d_A(X)\bigl(\tfrac bc+\tfrac cb\bigr)\ge 2\sum d_A(X)$, which is
$\sum d_A(X)\,(b-c)^2/(bc)\ge0$. In Lean this is one \lean{nlinarith} call with these products as
hints.

\paragraph{Equality.} Equality in the sum forces equality in each key lemma and
$d_A(X)(b-c)^2=d_B(X)(c-a)^2=d_C(X)(a-b)^2=0$. The equality case of the key lemma is the vanishing
of the square above (\lean{key\_scalar\_eq\_iff}), a linear relation between the weights for each
vertex. A combinatorial lemma (\lean{eq\_pos}) shows that these relations, together with the side
conditions (``if $w_1>0$ then \dots''), force the weights to be positive (\lean{eq\_pos} gives $w_0>0$; the
other two weights by the cyclic relabelling), hence every $d_i(X)>0$ and so $a=b=c$; then the linear relations give $w_0=w_1=w_2=\tfrac13$. The converse is a direct
computation for the equilateral triangle and its centroid. The scalar statement is
\lean{em\_eq\_scalar}; the geometric theorem identifies the weights $\tfrac13$ with Mathlib's centroid.

\paragraph{Barrow.} The proof avoids angles beyond the hypotheses. With
$x=|XA|$, $y=|XB|$, $z=|XC|$, $p=\langle XB,XC\rangle$, $q=\langle XC,XA\rangle$,
$\rho=\langle XA,XB\rangle$ (vectors from $X$):
(i) the angle-bisector theorem in squared form gives $|XU|^2(y+z)^2=2yz(yz+p)$
(\lean{bisector\_dist}), hence $|XU|\le\sqrt y\sqrt z\,\alpha$ with $\alpha^2=(yz+p)/(2yz)$, the
squared cosine of half of $\angle BXC$; (ii) since $X$ is strictly inside,
$w_0(A-X)+w_1(B-X)+w_2(C-X)=0$ with positive weights, which gives $xyz+xp+yq+z\rho\le0$ and the
vanishing of the $3\times3$ Gram determinant of the three vectors, and together these give
$\alpha^2+\beta^2+\gamma^2+2\alpha\beta\gamma=1$ for the three half-angle cosines; (iii) for
nonnegative $\alpha,\beta,\gamma$ with this relation, Wolstenholme's inequality
$\xi^2+\eta^2+\zeta^2\ge2(\alpha\eta\zeta+\beta\zeta\xi+\gamma\xi\eta)$ holds (a sum of squares), and with $\xi=\sqrt x$,
$\eta=\sqrt y$, $\zeta=\sqrt z$ it gives the result. The positivity of the two-dimensional Gram determinants
$\|B-X\|^2\|C-X\|^2-\langle B-X,C-X\rangle^2$ (needed for the bisector step) comes from the weight of
$A$ being nonzero (\lean{gram\_frame}).
\section{The Fermat--Torricelli point and the Napoleon configuration}\label{sec:fermat}
\subsection{Fermat--Torricelli (any dimension)}
These two theorems are in \lean{paperh\_mordell} and have no dimension hypothesis.
\begin{theorem}[Minimality]\label{thm:fermatmin}
Let $A,B,C,F$ be points of $P$ with $\angle AFB=\angle BFC=\angle CFA=2\pi/3$. Then for every
$X\in P$, $|FA|+|FB|+|FC|\le|XA|+|XB|+|XC|$.
\end{theorem}
\begin{theorem}[Existence and minimality]\label{thm:fermat}
Let $t$ be a triangle whose three angles $\angle BAC$, $\angle CBA$, $\angle ACB$ are all less than
$2\pi/3$. Then there is a point $F\in P$ with $\angle AFB=\angle BFC=\angle CFA=2\pi/3$, and for every
$X\in P$, $|FA|+|FB|+|FC|\le|XA|+|XB|+|XC|$.
\end{theorem}
In Theorem~\ref{thm:fermatmin} the points need not form a triangle: the angle hypotheses already
force $A,B,C\ne F$ (Mathlib's angle at a vertex coinciding with an endpoint is $\pi/2$). The angle
bound in Theorem~\ref{thm:fermat} is necessary for the existence of such an $F$; we do not formalise
the case of an angle $\ge 2\pi/3$, where the minimiser is the obtuse vertex, nor uniqueness of
the minimiser.

\paragraph{Proofs.} For minimality let $u_A,u_B,u_C$ be the unit vectors from $F$ to $A,B,C$. The
angle hypotheses give $\langle u_i,u_j\rangle=-\tfrac12$ for $i\ne j$, so $|u_A+u_B+u_C|^2=0$. By
Cauchy--Schwarz, $|XA|\ge\langle u_A,A-X\rangle=|FA|+\langle u_A,F-X\rangle$; summing the three
inequalities, the last terms cancel. For existence, the point is explicit
(\decl{ErdosMordell.fermat_exists}): with $g$ twice the area, $x_A=g+\sqrt3\langle B-A,C-A\rangle$
and cyclically, and $T=x_Bx_C+x_Ax_C+x_Ax_B$,
\[
F=A+\frac{x_Ax_C}{T}(B-A)+\frac{x_Ax_B}{T}(C-A),
\]
i.e.\ barycentric weights proportional to $1/x_A:1/x_B:1/x_C$. The angle condition $\angle A<2\pi/3$
is exactly $x_A>0$ (\lean{x\_pos}, using \lean{inner\_gt\_of\_angle\_lt}: $\angle(u,v)<2\pi/3$ implies
$\langle u,v\rangle>-\tfrac12|u||v|$), so the weights are positive ($F$ is interior; this is not stated in Lean), and the three angle
equalities reduce to polynomial identities in $b^2$, $c^2$, $k=\langle B-A,C-A\rangle$, $g$ and
$\sqrt3$ subject to $g^2=b^2c^2-k^2$ (\lean{fermat\_alg}, \lean{angle\_of\_alg}).

\subsection{The Fermat--Napoleon lines and Napoleon's theorem (oriented plane)}
The remaining theorems are in \lean{paperh\_fermat}, under the hypotheses
\lean{[Fact (finrank \(\mathbb{R}\) V = 2)]} and an orientation $o$ of $V$
(\lean{o : Orientation \(\mathbb{R}\) V (Fin 2)}). Let $\sigma\in\{1,-1\}$ be the orientation sign of
$t$ (\decl{FermatNapoleon.triSign}: $\sigma=1$ iff $\omega(B-A,C-A)>0$ for the area form $\omega$
of $o$). For points $X,Y$ the \emph{outer apex} on $XY$ is the image of $Y$ under the rotation about
$X$ by $-\sigma\pi/3$ (\decl{FermatNapoleon.outerApex}, with Mathlib's \lean{o.rotation}). Let $A'$,
$B'$, $C'$ be the outer apexes on $BC$, $CA$, $AB$ (\decl{FermatNapoleon.napoleonApex}) and $N_A$,
$N_B$, $N_C$ the centroids of the triangles $BCA'$, $CAB'$, $ABC'$
(\decl{FermatNapoleon.napoleonCenter}). The definitions use the orientation only to choose the
rotation direction; Theorem~\ref{thm:fn}(o) checks that the result does not depend on it in the
sense that matters.

\begin{theorem}[Fermat--Napoleon lines]\label{thm:fn}
In the setting above, for every triangle $t$:
\begin{enumerate}
\item[(o)] for each vertex $i$, the vertex $t_i$ and the apex on the opposite side lie strictly on
opposite sides of that side's line (Mathlib's \lean{SOppSide}); so the equilateral triangles are
erected outward;
\item[(i)] $|AA'|=|BB'|=|CC'|$;
\item[(ii)] there is a point lying on all three lines $AA'$, $BB'$, $CC'$;
\item[(iii)] if moreover $\angle BAC$, $\angle CBA$, $\angle ACB$ are all $<2\pi/3$, then every point
$F$ on the two lines $AA'$ and $BB'$ satisfies $\angle AFB=\angle BFC=\angle CFA=2\pi/3$ and
$|AA'|=|FA|+|FB|+|FC|$;
\item[(iv)] the (unoriented) angles between the vectors $A'-A$, $B'-B$, $C'-C$ are pairwise
$2\pi/3$, and the angle between $A'-A$ and $B-B'$ is $\pi/3$.
\end{enumerate}
\end{theorem}
\begin{theorem}[Napoleon; reproof]\label{thm:napoleon}
In the setting above, $|N_AN_B|=|N_BN_C|=|N_CN_A|$, and the centroid of $N_A,N_B,N_C$ is the
centroid $G$ of $t$.
\end{theorem}
Under the angle hypotheses of (iii), by (iii) and Theorem~\ref{thm:fermatmin} the common point of the lines is the Fermat point and
$|AA'|$ is the minimal value of $|XA|+|XB|+|XC|$; we do not state this combination as a separate
Lean theorem. Statement (iii) holds for any $F$ on the lines $AA'$ and $BB'$ (by symmetry for any two of the lines;
only this pair is formalised), hence for the common point of (ii); it uses the hypothesis $F\in BB'$ (the tampering suite removes it and the build fails) and
the angle hypothesis $\angle BAC<2\pi/3$ (likewise removed in a tampering test). In (iv), since lines
are unoriented, the angle between the lines $AA'$ and $BB'$ is $\pi/3$.

\paragraph{Proofs.} Let $J$ be Mathlib's right-angle rotation for $o$. The rotation by $-\sigma\pi/3$
is $v\mapsto\tfrac12v-\sigma\tfrac{\sqrt3}2Jv$ (\decl{FermatNapoleon.outerApex_vsub}). Every point
is written in the frame $p=B-A$, $q=C-A$ (\decl{FermatNapoleon.cramer}, using $\dim V=2$ and
$\omega(p,q)\neq0$), and $Jp$, $Jq$ are expressed in this frame through a scalar $\alpha>0$ with
$4\alpha^2(|p|^2|q|^2-\langle p,q\rangle^2)=3$ (\decl{FermatNapoleon.apex_frame}); this removes
$\sigma$ and $\omega$ from all later computations. With these coordinates, (o) is a sign computation
of $\omega$ (\decl{FermatNapoleon.sOppSide_of_areaForm}); (i) and (iv) follow from the Gram matrix of
$u_A=A'-A$, $u_B=B'-B$, $u_C=C'-C$, which has equal diagonal entries $M$ and off-diagonal entries
$-M/2$ (\decl{FermatNapoleon.apex_inner}). For (ii), with
$y_A=1+2\alpha\langle p,q\rangle$, $y_B=1+2\alpha(|p|^2-\langle p,q\rangle)$,
$y_C=1+2\alpha(|q|^2-\langle p,q\rangle)$ and $S=y_A+y_B+y_C>0$, the point $F$ with
$F-A=(y_A/S)\,u_A$ also satisfies $F-B=(y_B/S)\,u_B$ and $F-C=(y_C/S)\,u_C$
(\decl{FermatNapoleon.fermat_point_spec}; here $y_i=x_i/g$ in the notation of Section~\ref{sec:fermat}.1, and $u_A=A'-A$ already incorporates the factor $x_Bx_C$). For (iii), the angle conditions make
$y_A,y_B,y_C>0$ (\lean{weight\_pos}); if $F=A+r\,u_A=B+r'u_B$, taking inner products with $u_A$ and $u_B$
and using the Gram matrix gives $r=y_A/S$, so any $F$ on both lines is this point, and since $F-A$, $F-B$, $F-C$ are positive multiples of
$u_A,u_B,u_C$ the angles are those of (iv), i.e.\ $2\pi/3$; finally
$|FA|+|FB|+|FC|=\frac{y_A+y_B+y_C}{S}\,|u_A|=|AA'|$ because $|u_A|=|u_B|=|u_C|$. For Napoleon's theorem,
$N_B-N_A=\tfrac13(u_B-u_A)$ and cyclically (\decl{FermatNapoleon.napoleonCenter_vsub},
\decl{FermatNapoleon.apex_sum}), and the Gram matrix of the $u_i$ makes the three differences equal
in length; the centroid statement is $u_A+u_B+u_C=0$.

The argument for (ii)--(iii) is the classical one in coordinates; for the classical sources see
\cite{coxetergreitzer}. Napoleon's theorem was proved earlier in Lean in complex coordinates by
ourselves \cite{paperA}; the present proof is independent of that one and works in Mathlib's oriented
inner product plane.
\section{Verification}\label{sec:verify}
The checks follow our earlier papers \cite{paperA,paperG}; for a self-contained reading, they are:
\begin{center}\small
\begin{tabular}{@{}l p{11.2cm}@{}}
\toprule
gate & script that rejects forbidden tokens, builds, and checks that every required declaration uses only
the standard axioms\\
tampering suite & planted defects (hidden \lean{sorry}, extra axiom, fake axiom report, removed
hypothesis, false statement, \ldots) that the gate must reject\\
battery & per-theorem mutation of conclusions (each mutant must fail to compile; a heartbeat limit
counts as a rejection), a contradiction probe on the hypotheses, a satisfiability certificate,
hypothesis ablation, and the unused-variable linter\\
cross-machine build & the same sources built and gated on a second machine of a different
architecture\\
Layer 13 & blind restatement of the theorems from this paper's prose by a fresh writer, then a
two-way comparison with the library statements\\
Layer 14 & restatement of the main theorems in Mathlib-only terms, with explicit witness triangles\\
Layer 15 & statements written from classical sources by a writer who saw neither paper nor code, then
proved from the library\\
\bottomrule
\end{tabular}
\end{center}
The two machines are an x86-64 Linux machine and an arm64 machine. All builds use Lean v4.33.1 and
Mathlib at revision \texttt{0df444a360eaa60ab8c11dca51a86af692955474}.

\paragraph{Gates and tampering.} Each development has a gate script (\lean{gate.sh}) that rejects
forbidden tokens (\lean{sorry}, \lean{admit}, \lean{native\_decide}, \lean{axiom} declarations,
\lean{set\_option}, macros, notation, \lean{\#eval} and the like), builds the library, generates its
own \lean{\#print axioms} report for every declaration in its \lean{REQUIRED} list, and checks that
each depends only on \lean{propext}, \lean{Classical.choice} and \lean{Quot.sound}. The \lean{REQUIRED}
lists contain every definition, main theorem, bridge and auxiliary lemma. The results:
\begin{center}\small
\begin{tabular}{@{}l r l p{6.6cm}@{}}
\toprule
development & gate & tampering & removed hypotheses / false statements\\
\midrule
\lean{paperh\_euler} & PASS (60) & 9 of 9 rejected & weight-sum hypothesis of the two barycentric
distance lemmas; one Heron hypothesis of \lean{pedoe\_identity}; Euler's relation changed to
$R^2-3Rr$\\
\lean{paperh\_mordell} & PASS (48) & 7 of 7 rejected & ``$0\le w_i$'' weakened to a tautology in
\lean{erdos\_mordell}; $\angle CFA=2\pi/3$ dropped from \lean{fermat\_min}\\
\lean{paperh\_fermat} & PASS (47) & 9 of 9 rejected & $\angle BAC<2\pi/3$ and $F\in BB'$ dropped from
\lean{napoleon\_fermat}; $A\ne B$ dropped from \lean{ptolemy\_eq\_iff}; $\pi/3\mapsto\pi/4$ in
\lean{napoleon\_angle}\\
\bottomrule
\end{tabular}
\end{center}
In every suite the other five tests are a hidden \lean{sorry}, an axiom split over two lines, a fake
\lean{\#print axioms} report, a command-redefining macro and a wrapper axiom; each was rejected at the
expected stage (forbidden token, nonstandard axiom, or build failure). In \lean{paperh\_euler} the
main theorems have no hypothesis beyond the triangle, so the removed-hypothesis tests target the
lemmas through which they are proved. The gate counts in the table are those of the libraries
before the Layer-13--15 files were added. At the published commits, which add those files
(Section~\ref{sec:avail}), the gates print \lean{PASS (243 declarations, standard axioms only)}
(\lean{paperh\_euler}), \lean{PASS (144 \ldots)} (\lean{paperh\_mordell}) and \lean{PASS (171 \ldots)}
(\lean{paperh\_fermat}), and the three tampering suites again reject every planted defect.

\paragraph{Cross-machine builds.} \lean{paperh\_euler} and \lean{paperh\_fermat} were developed on
the x86-64 machine and rebuilt on the arm64 machine, \lean{paperh\_mordell} the other way round; the
gate printed \lean{PASS (60 declarations, standard axioms only)}, \lean{PASS (47 \ldots)} and
\lean{PASS (48 \ldots)} respectively on the second machine. For \lean{paperh\_mordell} the SHA-256
digest of the library files agrees between the two copies.

\paragraph{Fresh-clone receipts.} All three repositories were cloned from GitHub at the published
commits and built on both machines (x86-64 Linux and arm64 macOS, Lean v4.33.1, the pinned
Mathlib). On each machine the gates printed \lean{PASS (243 declarations, standard axioms only)},
\lean{PASS (144 \ldots)} and \lean{PASS (171 \ldots)} for \lean{paperh\_euler},
\lean{paperh\_mordell} and \lean{paperh\_fermat}; the tampering suites printed
\lean{ALL TAMPERING TESTS REJECTED}; and the SHA-256 digest of the sorted library \lean{.lean} file
hashes agreed between the machines (\texttt{66134350\ldots}, \texttt{922a0515\ldots} and
\texttt{02ca3ed9\ldots}). The same check was made at the first public commits, before the Layer-13--15
files (gates 60, 48 and 47; digests \texttt{d40a6efb\ldots}, \texttt{ca87ad23\ldots},
\texttt{3a5bd3d4\ldots}). The receipts are in the supporting material.

\paragraph{Statement-level battery.} The battery was run once on each development, on the machine
indicated:
\begin{center}\small
\begin{tabular}{@{}l l r r r r r@{}}
\toprule
development & machine & thms & mutants & rejected & survived & ablation\\
\midrule
\lean{paperh\_euler} & x86-64 & 56 & 228 & 225 (2 by timeout) & 3 & 29/0\\
\lean{paperh\_mordell} & arm64 & 47 & 162 & 161 & 1 & 211/0\\
\lean{paperh\_fermat} & x86-64 & 43 & 165 & 163 & 2 & 60/0\\
\bottomrule
\end{tabular}\\[2pt]
{\footnotesize ablation: hypotheses whose replacement by \lean{True} makes the existing proof fail/survive}
\end{center}
The ablation tests whether the supplied proof depends on a hypothesis, not whether the hypothesis is
logically necessary. The contradiction probe found no contradictory hypotheses (11, 42 and 20 declarations probed; 45, 5
and 23 have no hypotheses beyond the triangle or no hypotheses at all); the automatic satisfiability
certifier certified 7 of 56, 11 of 47 and 17 of 43 declarations, the rest being unresolved (not failures;
for \lean{paperh\_mordell} the satisfiability certifier rejected its own restatement of
\lean{em\_eq\_scalar} as ill-formed); the
unused-variable linter reported no warnings across the 14 theorem files (6, 4 and 4 files; the root import files and the definition-only \lean{Basic.lean} files have no theorems). All six survivors replace a statement by a
\emph{weaker true} statement, so their survival is expected and is not a defect of the statements:
\begin{itemize}
\item in \decl{EulerIneq.kooi_of_quartic}, $(4R+r)^2\mapsto(5R+r)^2$, which enlarges the upper bound
on $u$ (the mutant of the lemma compiles; the downstream proof of \lean{kooi}, which needs the exact
constant, breaks);
\item in \decl{EulerIneq.kooi_le_gerretsen}, the denominator $2(2R-r)\mapsto3(2R-r)$, which makes the
left side smaller;
\item in \decl{EulerIneq.hf_L_nonneg}, $2(ab+bc+ca)-(a^2+b^2+c^2)\ge0\mapsto
3(ab+bc+ca)-(a^2+b^2+c^2)\ge0$, implied by the original since $ab+bc+ca>0$;
\item in \lean{inner\_gt\_of\_angle\_lt} (present in both \lean{ErdosMordell} and
\lean{FermatNapoleon}, so counted in both batteries), $-\tfrac12|u||v|<\langle u,v\rangle\mapsto
-\tfrac22|u||v|<\langle u,v\rangle$, which is weaker and still true for nonzero $u,v$ with
$\angle(u,v)<2\pi/3<\pi$;
\item in \lean{weight\_pos}, $0<1+2\alpha m\mapsto0<2+2\alpha m$, which follows from the original.
\end{itemize}
Every audited theorem except \decl{FermatNapoleon.sOppSide_of_areaForm} (whose conclusion is an
\lean{SOppSide} predicate) had at least one conclusion mutant tested; most main theorems have no
hypotheses beyond the triangle and so are not probed for contradictions.

Finally the battery was run on fresh clones of the published commits, which include the
Layer-13--15 files, on both machines:
\begin{center}\small
\begin{tabular}{@{}l r r r r r r@{}}
\toprule
development & thms & mutants & rejected & survived & contradictions & sat.\ certified\\
\midrule
\lean{paperh\_euler} & 177 & 503 & 497 (2 by timeout) & 6 & 0 & 74 / 77\\
\lean{paperh\_mordell} & 101 & 314 & 310 & 4 & 0 & 36 / 39\\
\lean{paperh\_fermat} & 118 & 313 & 304 & 9 & 0 & 57 / 56\\
\bottomrule
\end{tabular}\\[2pt]
{\footnotesize sat.\ certified: x86-64 / arm64 (the certifier is model-assisted and its counts vary
slightly between runs; the other figures agree exactly between the machines)}
\end{center}
The hypothesis ablation found every tested hypothesis needed (33, 246 and 182 binders) and the
linter reported no warnings in the library files (9, 7 and 7 files). The contradiction probe found
no contradictory hypotheses; the satisfiability step marked 0, 4 and 2 of its own restatements as
ill-formed. The survivors are, on both machines, the same nineteen, and none is a defect of a
statement: six are the survivors classified above (\lean{kooi\_of\_quartic},
\lean{kooi\_le\_gerretsen}, \lean{hf\_L\_nonneg}, \lean{inner\_gt\_of\_angle\_lt} in both
\lean{ErdosMordell} and \lean{FermatNapoleon}, and \lean{weight\_pos}); the other thirteen are in the new
Layer-14 and Layer-15 files and replace a statement by a weaker true one. They are: the numerical
helper $\sqrt3<7/4\mapsto\sqrt3<8/4$; the conclusion $\angle ABC<2\pi/3\mapsto\angle ABC<\pi$ of
the helper \lean{l14\_angle\_lt} (in both \lean{ErdosMordell} and \lean{FermatNapoleon}), which derives
an angle bound from a nonnegative inner product and so also gives the weaker bound; a sign in the witness lemma \lean{l14T\_sign}, whose two branches
both equal $1$ after the mutation; two mutants each of the witness theorems \lean{l14w\_pedoe} and
\lean{l14w\_erdos\_mordell}, which change an equation fixing the witness data into a disequation or
flip a comparison inside a negated clause, leaving an existential statement that the same witness or
a different one still satisfies; and five mutants of \decl{FermatNapoleon.Classical.cov_F2d}, the
classical Fermat--Napoleon angle statement, whose conclusion is a disjunction
``angle $=\pi/3$ or angle $=2\pi/3$'' for each pair of lines: the library proves the second
disjunct, so mutating the first disjunct (or its vertex indices) leaves a true statement.

\paragraph{Statements in Mathlib's language and non-vacuity (Layer~14).} A file \lean{L14.lean} in
each development restates every main theorem with a type that mentions only Mathlib notions, our
definitions being unfolded: the area as $\tfrac12|AB|\,|AC|\sin A$, the semiperimeter as
$(a+b+c)/2$, the distance to a side line as Mathlib's \lean{Metric.infDist} to
\lean{line[\(\mathbb{R}\), Y, Z]}, the closed and open triangle as Mathlib's \lean{closedInterior} and
\lean{interior}, and the outer apexes as explicit rotations by $\mp\pi/3$ with the orientation sign
written out. Each restatement is proved from the library theorem. A check with fully qualified
pretty-printing found no constant of the three libraries in the elaborated types of the restatements
or witness theorems. Each witness theorem proves, on an explicit configuration in
\lean{EuclideanSpace \(\mathbb{R}\) (Fin 2)} or in $\mathbb{C}$, that every hypothesis holds and that the
conclusion is not degenerate:
\begin{itemize}
\item for the Euler family, the 3-4-5 triangle $A=(0,0)$, $B=(4,0)$, $C=(0,3)$ with $R=5/2$, $r=1$,
$s=6$, $\Delta=6$: $2r<R$ strictly, $|OI|^2=5/4$, the quartic takes the value $-1<0$, and both
Gerretsen bounds, Kooi's bound and the Weitzenb\"ock, Hadwiger--Finsler and isoperimetric
inequalities are strict; the Neuberg--Pedoe witness pairs this triangle with the non-similar
$(0,0),(1,0),(0,1)$ (strict, $48<50$), and Ptolemy's inequality is strict for the rectangle
$(0,0),(4,0),(0,3),(4,3)$ taken in non-cyclic order;
\item for Erdős--Mordell and the key lemma, the same triangle with $X=(1,1)$ in its interior (all
inequalities strict), and the equilateral triangle $(0,0),(2,0),(1,\sqrt3)$ with $X$ its centre for
the equality case; for Barrow, $A=(5,0)$, $B=(-3,4)$, $C=(-3,-4)$, $X=(0,0)$ with bisector feet
$(-3,0)$, $(1,-2)$, $(1,2)$ (strict); for Fermat--Torricelli, the 3-4-5 triangle (all angles
$<2\pi/3$, Fermat point distinct from the vertices) and, for minimality, $F=0$ with $A=(2,0)$,
$B=(-1,\sqrt3)$, $C=(-1,-\sqrt3)$;
\item for Ptolemy's equality, the rectangle $0,4,4+3i,3i$ (cyclic, diagonals crossing) and the
collinear points $0,1,2,3$; for the Fermat--Napoleon lines and Napoleon's theorem, $A=0$, $B=4$,
$C=3i$, with the apexes computed in $\mathbb{Q}(\sqrt3)$, pairwise distinct, the Fermat point on all
three lines and distinct from the vertices, and the Napoleon triangle non-degenerate.
\end{itemize}
The equality cases with an equilateral triangle or exact angles $2\pi/3$ need coordinates in
$\mathbb{Q}(\sqrt3)$; all other witnesses have integer coordinates. No restatement turned out false or
needed a hypothesis that the library theorem lacks.

\paragraph{Blind restatement from this paper (Layer~13).} A version of this paper with every Lean
name, listing, table of declarations and source link removed was given to a fresh AI writer with no
access to the code. From the prose alone it wrote 50 Lean statements for the 21 numbered theorems,
with its own definitions, compiled with only \lean{sorry} warnings; its session log shows that it
read only that text and its own files. For each pair, a comparison file proves the blind statement
from the library statement and conversely, using only unfolding, agreement lemmas between the two
sets of definitions and general Mathlib lemmas, never the library's own theorems. All 50 pairs agree.
Where this paper groups several claims in one theorem and the library proves them as one statement
(the three excentres, the two Gerretsen bounds, the three outward-apex claims, and the parts of
Theorems~\ref{thm:fn} and~\ref{thm:napoleon}), each blind statement follows from the library
statement and the blind statements of the group together imply it. The blind writer's eight
explicit collinear orders in Theorem~\ref{thm:ptolemyeq}(b) were shown equivalent to the library's
formulation for distinct collinear points. No difference between the prose and the Lean statements
was found. The comparison files, \lean{L13Blind.lean} in each library, compile with no \lean{sorry}
and only the standard axioms.

\paragraph{Comparison with the classical statements (Layer~15).} A second fresh writer, given only
textbook statements of twenty classical results (E1--E12: Euler, Euler's inequality, the excircle
relation, $OG$, $OH$, $IG$, the fundamental inequality, Gerretsen, Kooi, Weitzenb\"ock,
Hadwiger--Finsler, the isoperimetric inequality, Leibniz, Neuberg--Pedoe; P1--P2: Ptolemy's
inequality and its equality case; M1--M2: Erdős--Mordell and Barrow; F1--F3: Fermat--Torricelli,
the Fermat--Napoleon lines, Napoleon) and neither this paper nor the code, wrote 33 Lean statements
with its own definitions. Comparison files prove 32 of them from the libraries. For the classical
form of the Fermat--Napoleon statements, where the apexes are any points forming outward equilateral
triangles, and for the classical Ptolemy equality (``in order on a circle or a line''), new bridge
lemmas were proved: a point forming an outward equilateral triangle on a side is the library's apex,
and a point seeing the sides at $2\pi/3$ is unique (so it is the common point of the lines). The
remaining statement, P1, is Ptolemy's inequality as the writer typed it: because the statement does
not mention the vector space, Lean bound it to an arbitrary metric space, where it is false (the
comparison file proves a counterexample in the plane with the maximum metric). Its intended
Euclidean form is exactly Theorem~\ref{thm:ptolemy}. These files are \lean{L15Classical.lean} in each
library.

\paragraph{Hand check and external review.} The author checked the numbered theorems against the
Lean statements by hand. The draft was reviewed by three AI models other than the one that wrote it,
and all their findings were addressed before publication; one of them, an independent referee
review, asked that the three developments be committed and their copies synchronised across the two
machines before the final figures were taken, which the published commits and the receipts above
now record.
\section{Related work and prior art}\label{sec:related}
\paragraph{Euler's relation and inequality.} Weiyi Wang (\lean{wwylele}) proved, in a public Lean~4
repository on Mathlib's \lean{Affine.Triangle} API (file first committed 27 January 2026)
\cite{wwylele}, $|IO|^2=R(R-2r)$ (\lean{dist\_incenter\_circumcenter}), a signed excentre version,
and $2r\le R$ (\lean{two\_mul\_inradius\_le\_circumradius}). The same author's Mathlib pull request
\#44230 \cite{pr44230} (opened 26 September 2026; open, labelled as blocked and with merge conflicts
when we checked) contains \lean{dist\_incenter\_circumcenter\_sq},
\lean{two\_mul\_inradius\_le\_circumradius} and
\lean{dist\_excenter\_singleton\_circumcenter\_sq}, which are, up to ring normalisation, our
Theorems~\ref{thm:euler} and~\ref{thm:eulerex} and the inequality of Theorem~\ref{thm:eulerineq}.
The repository \lean{facebookresearch/atlas-lean} \cite{atlaseuler} (5 October 2026) proves
$|OI|^2=R(R-2r)\wedge R\ge2r$ in \lean{EuclideanSpace \(\mathbb{R}\) (Fin 2)}, with the circumcircle
and incircle given by explicit data (feet of perpendiculars and positive weights). We found no
formal proof of the equality case of Euler's inequality in any of these (no \lean{equilateral} or
\lean{eq\_iff} statement). Our proofs were written without using these developments; they share the
barycentric route through Mathlib's excentre weights, which is the natural one in this API. If
\#44230 is merged, Theorems~\ref{thm:euler}, \ref{thm:eulerex} and the inequality of
Theorem~\ref{thm:eulerineq} become Mathlib lemmas.

\paragraph{Ptolemy.} Ptolemy's inequality in any real inner product space,
\mdecl{EuclideanGeometry.mul_dist_le_mul_dist_add_mul_dist}, was added to Lean~3 mathlib by
Kudryashov in June 2022 and ported to Mathlib; it is identical to our Theorem~\ref{thm:ptolemy},
which merely restates it. \lean{lean-genius} has another proof (24 June 2026) \cite{lgptolemy}. The
equality for concyclic points (Ptolemy's theorem) is in Mathlib
(\mdecl{EuclideanGeometry.mul_dist_add_mul_dist_eq_mul_dist_of_cospherical}), HOL~Light, Isabelle,
Rocq, Metamath and Mizar (Wiedijk's list, theorem~95 \cite{wiedijk100}). We found no formal proof of
the converse, i.e.\ of the characterisation in Theorem~\ref{thm:ptolemyeq}.

\paragraph{Erdős--Mordell.} Two sorry-free Lean~4 proofs of the inequality predate ours, both for a
point in the topological interior of a non-degenerate triangle in the plane: \lean{erdos\_898} in
\lean{plby/lean-proofs} (earliest copy 7 May 2026; any two-dimensional inner product space, with
distances to lines through \lean{orthogonalProjection}) \cite{plby}, produced with AI assistance and
linked as the formal proof of Erdős problem~898 from Google DeepMind's \lean{formal-conjectures}
\cite{formalconj}; and \lean{erdos\_mordell} in \lean{lean-genius} (completed 19 June 2026;
\lean{EuclideanSpace \(\mathbb{R}\) (Fin 2)}), with the key lemma
\lean{key\_inequality\_A/B/C} in the same weighted form as our Theorem~\ref{thm:key}
\cite{lgmordell}. The equality case appears in \lean{lean-genius} as an \lean{axiom}
(\lean{erdos\_mordell\_equality}, January 2026), in a file that also contains \lean{sorry}
\cite{lg898}; we found no proof of it. A variant with constant $3/2$ in arbitrary normed planes is in
the Principia \lean{math-solutions} collection (26 October 2026) \cite{principia}; it is a different
statement. We found no formal proof of Barrow's inequality.

\paragraph{Weitzenb\"ock, Hadwiger--Finsler, isoperimetric, Leibniz.} Weitzenb\"ock's inequality
with its equality case (IMO 1961 Problem 2) is proved in Compfiles (\lean{Imo1961P2}, December 2025)
\cite{compfiles} and in \lean{lean-genius} (June 2026) \cite{lgweitz}, and Hadwiger--Finsler with
its equality case in \lean{lean-genius} (June 2026) \cite{lghf}. All three are algebraic: the
``area'' is a real number $T$ assumed to satisfy Heron's relation $16T^2=(a+b+c)(-a+b+c)\cdots$
(with the triangle inequalities as further hypotheses for Hadwiger--Finsler). Our
Theorems~\ref{thm:weitz} and~\ref{thm:hf} are about the area of an actual \lean{Affine.Triangle},
which requires the Heron bridge; the algebraic core is the same, so the novelty here is limited to
the geometric statement. (The ``Finsler--Hadwiger theorem'' in \lean{atlas-lean} is a different
result about squares erected on two sides.) For the triangle isoperimetric inequality and the
Leibniz minimum we found only the general isoperimetric theorem (HOL~Light) and unrelated
variance-type identities; we did not exhaustively search Mathlib's analysis files for a Leibniz-type
identity.

\paragraph{Gerretsen, Kooi, $OG$, $OH$, $IG$, Neuberg--Pedoe, Fermat--Torricelli.} We found no formal
proof of any of these. Mathlib has the distance from a vertex to a point of the Euler line
(\lean{dist\_lineMap\_circumcenter\_centroid\_sq}, merged in October 2026 after our pinned revision)
and the nine-point-circle centre, from which $OH^2$ follows quickly, but not $OG^2$, $OH^2$ or
$IG^2$ themselves. The only Lean hit for the Fermat point is a computable bound in a Steiner-ratio
certificate project, not a theorem. Our earlier Blundon development \cite{paperD} proves Blundon's
inequalities, from which the fundamental quartic is taken.

\paragraph{Napoleon.} Napoleon's theorem is formalised in Isabelle (AFP entry by Ramos, Hulak and de
Queiroz, August 2026, in complex coordinates with both orientations) \cite{afpnapoleon}; in Lean by
Tanr\i{}kolu (complex coordinates, the algebraic equilateral criterion) \cite{tanrikolu}, in the
\lean{icml-26-lean-challenges} benchmark (June 2026, distance equalities for both Napoleon triangles)
\cite{shadowbench}, in \lean{atlas-lean} (October 2026, \lean{EuclideanSpace \(\mathbb{R}\) (Fin 2)},
apexes given by distance equalities and signed areas) \cite{atlasnapoleon}; and in our own Paper~A
\cite{paperA} (complex coordinates, both Napoleon triangles). Our Theorem~\ref{thm:napoleon} is a
further independent proof, in Mathlib's oriented plane with apexes defined by \lean{o.rotation},
which adds the centroid statement. Guilhot's Coq development of oriented angles includes a file on
Napoleon's theorem \cite{guilhot}. We found no formal proof of the Fermat--Napoleon statements of
Theorem~\ref{thm:fn}(i)--(iv) (the equal lengths $AA'=BB'=CC'$, the concurrence, the Fermat point on
the lines with $AA'=FA+FB+FC$, and the angles between the lines); the \lean{atlas-lean} Napoleon file
does not contain them. This claim rests on a narrower search than the rest of this section: only the
AFP entry page, the \lean{atlas-lean} file and a GitHub code search were checked for it, and the
GeoCoq, HOL~Light and AFP theory sources were not searched for the concurrence.

\paragraph{Automated reasoning.} Many of these inequalities have been proved by automated methods
outside interactive provers, for instance by quantifier elimination and with dynamic-geometry tools
(see \cite{kovacs2017} for Euler's inequality in GeoGebra). These are not kernel-checked proofs about
library objects; we did not survey this literature systematically.

\paragraph{How the search was done, and its gaps.} The sweep (28 October 2026) used the GitHub API
(code, pull-request and repository search in \lean{leanprover-community/mathlib4} and globally, all
214 pull requests labelled \lean{t-euclidean-geometry}), full tree listings of \lean{lean-genius},
\lean{atlas-lean}, Compfiles, \lean{formal-conjectures}, \lean{plby/lean-proofs} and
\lean{wwylele/Feuerbach}, the Lean~3 mathlib and miniF2F trees, HOL~Light, GeoCoq, Metamath
(\lean{set.mm}), the Mizar EUCLID articles, the AFP geometry topic (33 entries), Wiedijk's list and the
arXiv API. Every file classified as prior art was downloaded and read; none was compiled, so
``sorry-free'' means that a text search found no \lean{sorry}, \lean{admit} or \lean{axiom}. The
Napoleon prior art comes from the search for our Paper~A \cite{paperA}, re-checked on 28 October 2026
by fetching the AFP entry page and the \lean{atlas-lean} file. Gaps: general web search was not
usable during the sweep, so Google Scholar, the Lean Zulip, mailing lists and conference proceedings
(ITP, CICM, ADG, ThEdu) were not searched; GitHub code search covers only default branches and caps
its results, and two rate-limited queries were not re-run; AFP theory sources outside the geometry
topic, the full HOL~Light \lean{Multivariate} tree, Mizar articles other than EUCLID, and
fast-growing AI-proof collections were covered only partially. Our statements of novelty are
subject to these gaps.

\paragraph{Classical sources.} For the inequalities we follow Bottema et al.\ \cite{bottema}; the
original sources are Weitzenb\"ock \cite{weitzenboeck}, Finsler and Hadwiger \cite{finslerhadwiger},
Pedoe \cite{pedoe1941,pedoe1942}, Blundon \cite{blundon1965}, Erdős's problem \cite{erdos1935} with
the solution by Mordell and Barrow \cite{mordellbarrow1937}, and Kazarinoff's proof of the key lemma
\cite{kazarinoff1957}. Gerretsen's and Kooi's inequalities are cited, as is usual, through
\cite{bottema}; we did not obtain the original papers. For the Fermat point and Napoleon's theorem
see Coxeter and Greitzer \cite{coxetergreitzer}.

\section{Limitations}\label{sec:limits}
\begin{enumerate}
\item \textbf{Reproofs.} Theorems~\ref{thm:euler}, \ref{thm:eulerex}, the inequality of
\ref{thm:eulerineq}, \ref{thm:ptolemy}, \ref{thm:key}, the inequality of \ref{thm:em} and
\ref{thm:napoleon} were formalised before us (Section~\ref{sec:related}); we claim only independent
proofs, and for Erdős--Mordell the mild extension to the closed triangle in any dimension.
\item \textbf{Equality cases not formalised.} Kooi's and Gerretsen's inequalities, Barrow's
inequality and the Fermat--Torricelli minimality are stated without equality cases; uniqueness of the
Fermat point and the case of an angle $\ge2\pi/3$ are not formalised.
\item \textbf{Dimension and orientation.} Theorems~\ref{thm:fn} and~\ref{thm:napoleon} assume
$\dim V=2$ and an orientation; all other theorems hold in any dimension.
\item \textbf{One labelling, one statement form.} Where a statement has three cyclic forms (key lemma,
excircles) all are stated; otherwise no relabelling lemmas are proved. Barrow's bisector feet are
given by hypotheses (betweenness and equal angles), not constructed.
\item \textbf{Our own definitions.} The area, the semiperimeter, the orientation sign, the outer apexes
and the Napoleon centres are our definitions; the area is linked to Mathlib by the Heron bridge, and
\ref{thm:fn}(o) checks that the apexes are outward. Gerretsen and Kooi are stated with our $s$.
\item \textbf{All-AI checking.} The proofs, numerical checks, battery and prior-art search were done by
AI agents; no human has checked the proofs line by line, which is the kernel's job.
\item \textbf{Prior art.} The search had the gaps listed in Section~\ref{sec:related}.
\end{enumerate}
\section{Data availability}\label{sec:avail}
The proofs are in three Lean~4 developments on Lean v4.33.1 with Mathlib pinned at revision
\texttt{0df444a3} (in \lean{lake-manifest.json}). Line counts are of the library \lean{.lean} files,
including the root import file, at the published commits.
\begin{center}\small
\begin{tabular}{lp{8.4cm}rr}
\toprule
repository & library files & files & lines\\
\midrule
\lean{paperh\_euler} & \lean{Basic} (32), \lean{Bridges} (330), \lean{Main} (402), \lean{Pedoe} (143), \lean{Ptolemy} (24), \lean{Kooi} (63), \lean{Extremal} (158), \lean{L14} (330), \lean{L13Blind} (606), \lean{L15Classical} (328) & 11 & 2{,}426\\
\lean{paperh\_mordell} & \lean{Basic} (116), \lean{Mordell} (431), \lean{Fermat} (306), \lean{Barrow} (328), \lean{L14} (381), \lean{L13Blind} (305), \lean{L15Classical} (167) & 8 & 2{,}041\\
\lean{paperh\_fermat} & \lean{Basic} (44), \lean{Lemmas} (81), \lean{Bridges} (346), \lean{Ptolemy} (352), \lean{Main} (274), \lean{L14} (381), \lean{L13Blind} (461), \lean{L15Classical} (472) & 9 & 2{,}419\\
\bottomrule
\end{tabular}
\end{center}
The repositories are \url{https://github.com/jpbald93/paperh_euler} at commit
\texttt{0c95102c792b8d47d2720f4e9fbc39aa5f21a424},
\url{https://github.com/jpbald93/paperh_mordell} at commit
\texttt{7b2270146195cf2a4c589e403143dba0abd0f6a7} and
\url{https://github.com/jpbald93/paperh_fermat} at commit
\texttt{8f680c86b068196a4d27201209c856a5d038aaa4}. In each, the first commit is the library as
described in Sections~\ref{sec:euler}--\ref{sec:fermat}, published as a single commit (the
development history is not included); the next two add the Layer-14 and the Layer-13/15 files,
without changing any earlier file. Every \texttt{File:line} reference in this paper links to the
published commit. Supporting material (the prior-art sweep, the builders' reports, the battery, gate,
tampering, cross-build and fresh-clone reports, and the Layer-13--15 reports) is in the folder
\lean{PaperH-triangle-inequalities} of the repository
\url{https://github.com/jpbald93/LEAN-verification}. Every name in \lean{\textbackslash decl} font
(qualified by its library namespace in Table~\ref{tab:thms}) is in the \lean{REQUIRED} list of the
repository that contains it at the published commit; a script checked all 88 distinct names (91
occurrences) against the published \lean{REQUIRED} lists, and all 33 \texttt{File:line} references
(and the further line numbers listed with them) against the declarations at those lines.

\section*{Acknowledgements and use of AI}
Large-language-model AI assistants were used to write the Lean proofs, to run the prior-art search,
the numerical checks and the verification battery, and to draft this paper; parts of the checking were
done by a different model from the one that wrote the proofs. The author directed the work and
reviewed the paper. Every proof is checked by the Lean kernel, and the author takes responsibility for
all content.
\begin{thebibliography}{99}
\bibitem{lean4}
L.~de Moura and S.~Ullrich,
\emph{The Lean 4 theorem prover and programming language},
in: Automated Deduction -- CADE-28, Lecture Notes in Computer Science \textbf{12699}, Springer, 2021, 625--635,
\url{https://doi.org/10.1007/978-3-030-79876-5_37}.

\bibitem{mathlib}
The mathlib Community,
\emph{The Lean mathematical library},
in: Proceedings of the 9th ACM SIGPLAN International Conference on Certified Programs and Proofs (CPP '20), ACM, 2020, 367--381,
\url{https://doi.org/10.1145/3372885.3373824}.

\bibitem{bottema}
O.~Bottema, R.~\v{Z}.~\DJ{}or\dj{}evi\'c, R.~R.~Jani\'c, D.~S.~Mitrinovi\'c and P.~M.~Vasi\'c,
\emph{Geometric Inequalities}, Wolters-Noordhoff, Groningen, 1969.

\bibitem{weitzenboeck}
R.~Weitzenb\"ock, \"Uber eine Ungleichung in der Dreiecksgeometrie, \emph{Mathematische Zeitschrift}
\textbf{5} (1919), 137--146, \url{https://doi.org/10.1007/BF01203160}.

\bibitem{finslerhadwiger}
P.~Finsler and H.~Hadwiger, Einige Relationen im Dreieck, \emph{Commentarii Mathematici Helvetici}
\textbf{10} (1937), 316--326, \url{https://doi.org/10.1007/BF01214300}.

\bibitem{erdos1935}
P.~Erd\H{o}s, Problem 3740, in: Problems for Solution: 3739--3743, \emph{The American Mathematical
Monthly} \textbf{42} (1935), 396, \url{https://doi.org/10.2307/2301373}.

\bibitem{mordellbarrow1937}
P.~Erd\H{o}s, L.~J.~Mordell and D.~F.~Barrow, Solution of Problem 3740, \emph{The American
Mathematical Monthly} \textbf{44} (1937), 252, \url{https://doi.org/10.2307/2300713}.

\bibitem{kazarinoff1957}
D.~K.~Kazarinoff, A simple proof of the Erd\H{o}s--Mordell inequality for triangles,
\emph{Michigan Mathematical Journal} \textbf{4} (1957), 97--98,
\url{https://doi.org/10.1307/mmj/1028988998}.

\bibitem{pedoe1941}
D.~Pedoe, An inequality connecting any two triangles (Note 1562), \emph{The Mathematical Gazette}
\textbf{25} (1941), 310--311, \url{https://doi.org/10.2307/3606570}.

\bibitem{pedoe1942}
D.~Pedoe, An inequality for two triangles, \emph{Mathematical Proceedings of the Cambridge
Philosophical Society} \textbf{38} (1942), 397--398, \url{https://doi.org/10.1017/S030500410002209X}.

\bibitem{blundon1965}
W.~J.~Blundon, Inequalities associated with the triangle, \emph{Canadian Mathematical Bulletin}
\textbf{8} (1965), 615--626, \url{https://doi.org/10.4153/CMB-1965-044-9}.

\bibitem{coxetergreitzer}
H.~S.~M.~Coxeter and S.~L.~Greitzer,
\emph{Geometry Revisited},
Mathematical Association of America, 1967.

\bibitem{wwylele}
W.~Wang (GitHub user \texttt{wwylele}), \emph{Feuerbach}, GitHub repository, file
\path{Feuerbach/Basic.lean} (first committed 27 January 2026),
\url{https://github.com/wwylele/Feuerbach}, accessed 28 October 2026.

\bibitem{pr44230}
W.~Wang, \emph{feat(Geometry/Euclidean): prove Euler's theorem in geometry}, mathlib4 pull request \#44230, opened 26 September 2026 (open at the time of writing),
\url{https://github.com/leanprover-community/mathlib4/pull/44230}, accessed 28 October 2026.

\bibitem{atlaseuler}
\emph{atlas-lean} (Meta / facebookresearch), file \path{MathlibExt/Geometry/Euclidean/EulerOI.lean}
(5 October 2026),
\url{https://github.com/facebookresearch/atlas-lean/blob/main/MathlibExt/Geometry/Euclidean/EulerOI.lean},
accessed 28 October 2026.

\bibitem{plby}
GitHub user \texttt{plby}, \emph{lean-proofs}, file \path{src/latest/ErdosProblems/Erdos898.lean}
(earliest copy 7 May 2026),
\url{https://github.com/plby/lean-proofs/blob/main/src/latest/ErdosProblems/Erdos898.lean},
accessed 28 October 2026.

\bibitem{formalconj}
Google DeepMind, \emph{formal-conjectures}, file \path{FormalConjectures/ErdosProblems/898.lean}
(4 August 2026),
\url{https://github.com/google-deepmind/formal-conjectures/blob/main/FormalConjectures/ErdosProblems/898.lean},
accessed 28 October 2026.

\bibitem{lgmordell}
GitHub user \texttt{rjwalters}, \emph{lean-genius}, file
\path{proofs/Proofs/ErdosMordellInequalityOQ01.lean} (completed 19 June 2026),
\url{https://github.com/rjwalters/lean-genius/blob/main/proofs/Proofs/ErdosMordellInequalityOQ01.lean},
accessed 28 October 2026.

\bibitem{lg898}
GitHub user \texttt{rjwalters}, \emph{lean-genius}, file \path{proofs/Proofs/Erdos898Problem.lean}
(15 January 2026),
\url{https://github.com/rjwalters/lean-genius/blob/main/proofs/Proofs/Erdos898Problem.lean},
accessed 28 October 2026.

\bibitem{lgweitz}
GitHub user \texttt{rjwalters}, \emph{lean-genius}, file
\path{proofs/Proofs/WeitzenbockInequalityOQ01.lean} (16 June 2026),
\url{https://github.com/rjwalters/lean-genius/blob/main/proofs/Proofs/WeitzenbockInequalityOQ01.lean},
accessed 28 October 2026.

\bibitem{lghf}
GitHub user \texttt{rjwalters}, \emph{lean-genius}, file
\path{proofs/Proofs/WeitzenbockInequalityOQ01OQ01.lean} (23 June 2026),
\url{https://github.com/rjwalters/lean-genius/blob/main/proofs/Proofs/WeitzenbockInequalityOQ01OQ01.lean},
accessed 28 October 2026.

\bibitem{lgptolemy}
GitHub user \texttt{rjwalters}, \emph{lean-genius}, file
\path{proofs/Proofs/PtolemyInequalityInnerProduct.lean} (24 June 2026),
\url{https://github.com/rjwalters/lean-genius/blob/main/proofs/Proofs/PtolemyInequalityInnerProduct.lean},
accessed 28 October 2026.

\bibitem{compfiles}
D.~Renshaw et al., \emph{Compfiles}, file \path{Compfiles/Imo1961P2.lean} (30 December 2025),
\url{https://github.com/dwrensha/compfiles/blob/main/Compfiles/Imo1961P2.lean}, accessed 28 October 2026.

\bibitem{principia}
Principia, \emph{math-solutions}, file \path{mathdb/Principia/MathDB/Problem341671.lean}
(26 October 2026),
\url{https://github.com/PrincipiaMathOrg/math-solutions/blob/main/mathdb/Principia/MathDB/Problem341671.lean},
accessed 28 October 2026.

\bibitem{wiedijk100}
F.~Wiedijk, \emph{Formalizing 100 Theorems}, \url{https://www.cs.ru.nl/~freek/100/}, accessed 28 October 2026.

\bibitem{afpnapoleon}
A.~F.~Ramos, D.~B.~Hulak and R.~J.~G.~B.~de Queiroz, Napoleon's Theorem, \emph{Archive of Formal
Proofs}, August 2026, \url{https://isa-afp.org/entries/Napoleon.html}, accessed 28 October 2026.

\bibitem{tanrikolu}
I.~Tanr\i{}kolu, \emph{Napoleon-s-Theorem-in-Lean4}, file \path{napoleantheorem.lean}, commit
\texttt{0e6eb11}, \url{https://github.com/imran-tanrikolu/Napoleon-s-Theorem-in-Lean4}, accessed
28 October 2026.

\bibitem{shadowbench}
\emph{icml-26-lean-challenges} (ShadowBench), file
\path{autformalization/problems/geometry/L2/geo_gen_L2_006/ShadowBench/Source/Main.lean}, commit
\texttt{7541fcec} (12 June 2026), \url{https://github.com/epfl-lara/icml-26-lean-challenges}, accessed
28 October 2026.

\bibitem{atlasnapoleon}
\emph{atlas-lean} (Meta / facebookresearch), file
\path{MathlibExt/Geometry/Euclidean/NapoleonTheorem.lean}, pull request \#63, commit
\texttt{27a38cfd} (5 October 2026),
\url{https://github.com/facebookresearch/atlas-lean/blob/main/MathlibExt/Geometry/Euclidean/NapoleonTheorem.lean},
accessed 28 October 2026.

\bibitem{guilhot}
F.~Guilhot, \emph{Formalization of the oriented angles theory} (Coq contribution \texttt{Angles}, 2002;
file \path{point_napoleon.v}), \url{https://github.com/rocq-archive/angles}, accessed 28 October 2026.

\bibitem{kovacs2017}
Z.~Kov\'acs, R.~Vajda and A.~Montag, On Euler's inequality and automated reasoning with dynamic
geometry, arXiv:1708.02993, 2017, \url{https://arxiv.org/abs/1708.02993}.

\bibitem{paperA}
J.~Bald,
\emph{Classical plane geometry by coordinate certificates: Desargues, Monge, Simson, Napoleon and Lami in Lean~4},
Zenodo, 2026. \url{https://doi.org/10.5281/zenodo.23239930}

\bibitem{paperD}
J.~Bald,
\emph{Circle configurations in classical plane geometry, formalized in Lean~4: Clifford's chain, Miquel's five circles, Lester, Sawayama--Th\'ebault and Kosnita},
Zenodo, 2026. \url{https://doi.org/10.5281/zenodo.23378578}

\bibitem{paperG}
J.~Bald,
\emph{Triangle centres in Lean~4: the Nagel line family and the symmedian and Brocard points},
Zenodo, 2026. \url{https://doi.org/10.5281/zenodo.23741528}
\end{thebibliography}
\end{document}
